\section{Implications and remaining questions}\label{sec:discussion}

The three-variable results give a precise reason to strengthen a local
relaxation. The five edge contacts of the new family cannot be matched by
the allowed affine factors in a disjoint-support certificate. This failure
persists for a strictly positive quadratic and is witnessed by rational
moments that satisfy the specified recent SOC strengthening. Enforcing
the family does not require a larger hierarchy of moments: one order-five
block and six nonnegative scalars suffice for each orientation. The exact
tetrahedral hull remains an alternative when complete local convexification
is needed.

The closedness, rounding, and sign results narrow the remaining
three-variable question. A missing extreme ray must have all three square
coefficients positive and a positive product of mixed coefficients.
Cube symmetries then reduce its sign pattern, and the contact arguments
exclude several zero configurations. \Cref{three:edge-classification}
settles every sign pattern when all vertex values are positive: an
extreme ray with positive square coefficients is an affine square or a
symmetry copy of the strict family. Rays with a vertex zero still require
classification beyond the boundary results proved here. The precise
completeness question is stated in \cref{sec:three}; none of the validity,
enforcement, or representability results depends on a positive answer.
The sampled moment and contact studies provide numerical evidence for
the broader equality, rather than a proof.

The dimension threshold sets a different limit. In four or more variables,
no finite collection of affine semidefinite constraints and auxiliary
variables can represent the full quadratic moment hull, or its positive
diagonal enlargement. The obstruction is local to a vertex and extends
to other polytopes with a simple vertex. It does not prevent exact
separation, convergent hierarchies, or finite formulations for particular
objective classes. The proper-face result also shows that the obstruction
can disappear after restricting to the zero set of any nonzero supporting
quadratic in four variables.

For sparse models, local dimension alone is insufficient. Components with
at most three positive vertices can be combined through binary boundary
states to give a finite lift. A $K_4$ minor among positive vertices rules
out any finite lift. The intervening graph class remains open, with the
four-vertex path and the three-leaf star as the first cases identified
here. The available finite construction can also be large: its size
depends on the torso obtained by completing component boundaries, rather
than only on the original graph width.

The computational evidence supports a limited practical conclusion.
Selected family blocks can recover exact local-hull bounds on constructed
instances designed around the missing inequalities. The tested natural
instances give much less room for improvement, and the archived timings
do not establish which formulation is fastest under matched selection
and computing conditions. A solver implementation should therefore
measure both the bound gained and the cost of each selected block.
The exact separation and gain bounds in the appendices provide tools
for that comparison, provided their arithmetic and locality hypotheses
are respected.
